\documentclass[10pt,a4paper]{amsart}
\usepackage[margin=19mm]{geometry}
\usepackage{amsmath,amssymb,mathtools}
\usepackage[scr=rsfs,cal=euler]{mathalfa}
\usepackage{xfrac}
\usepackage[unicode,hidelinks,bookmarks=false]{hyperref}
\newcommand{\ie}{i.e.\ }
\newcommand{\cf}{cf.\ }
\newcommand{\resp}{respectively\ }
\newcommand{\Subsection}{\subsection}
\providecommand{\nobreakdash}{\nobreak\hbox{-}}
\setlength{\emergencystretch}{2em}
\multlinegap0pt
\usepackage{amssymb}
\usepackage{tikz}
\usepackage{mathtools}
\usepackage{xcolor}
\newtheorem{theorem}{Theorem}
\title{Sasaki-Einstein metrics on connected sums of $S^3\times S^4$}
\author{Dasol Jeong, In-Kyun Kim, Jihun Park, Joonyeong Won}
\date{Source: arXiv:2608.29695v1 (CC BY 4.0)}
\begin{document}
\maketitle

\noindent\emph{Source and licence.} Sasaki-Einstein metrics on connected sums of $S^3\times S^4$. Version arXiv:2608.29695v1 (CC BY 4.0), distributed by arXiv under the Creative Commons Attribution 4.0 International licence, http://creativecommons.org/licenses/by/4.0/, as recorded in the arXiv licence metadata for this exact version and shown on the abstract page https://arxiv.org/abs/2608.29695v1. Full licence notice: \texttt{licenses/arxiv-2608.29695-CC-BY-4.0.txt}.\par\smallskip

\noindent\textit{Editorial scope.} Counted unit: the article's theorem theorem:h (source0.tex lines 707--712). The article's complete original proof of it is delivered with it (source0.tex lines 713--1071, one \texttt{proof} environment of 359 lines). No mathematics is written, repaired or re-derived: the statement and the proof are the article's own lines, in the article's own notation, and no retained line is substituted or re-flowed.  Retained uncounted as the source-local context the proof needs: lines 446--449.  Nothing was omitted from any retained span, so the package carries no omission record.  The retained lines cite 2 works, whose entries are transcribed verbatim from the article's own bibliography source in the e-print and delivered with short editorial labels recorded in the item's reference mapping.  No retained line contains a figure environment, an image-inclusion command or drawing code, so no image is delivered and the package declares no asset dependency.  Editorial formatting only, in addition: an AMS \texttt{amsart} preamble that restates the article's own declarations and macros used by the retained lines, namely the environments \texttt{theorem} on separate counters, together with the standard packages those lines need.  The retained environments print in this document as Theorem 1 in order of appearance, whereas the article prints the same items with the numbering of its own full document.  The complete native source of the article as distributed in the arXiv e-print for this exact version is bundled unchanged as \texttt{sources/208880/source0.tex}.
\begin{theorem}[{\cite[Corollary~4.14]{Z21}}]\label{theorem:G-delta}
Let $(W, \Delta_W)$ be a  log $\mathbb{Q}$-Fano variety with a reductive algebraic subgroup~$\mathfrak{G}$ of $\mathrm{Aut}(W, \Delta_W)$.
A log $\mathbb{Q}$-Fano variety $(W, \Delta_W)$ is K-polystable  if $\delta_\mathfrak{G}(W, \Delta_W)>1$.
\end{theorem}

\begin{theorem}\label{theorem:h}
For each integer $n\geq 4$,  the hypersurface $Y$ of degree $12n$  defined by 
\[x_0^{n}+x_1^{n}+x_3x_2^4+x_1x_3^3+x_4^2=0\]
in $\mathbb{P}(12,12,2n+1,4n-4,6n)$
is K-polystable.
\end{theorem}

\begin{proof}
First, suppose that $n\not\equiv 1$ (mod $3$). 

The hypersurface $Y$ is a double cover of $\mathbb{P}(12,12,2n+1,4n-4)$ branched along the surface
\[
x_0^{n}+x_1^{n}+x_3x_2^4+x_1x_3^3=0.
\]
After dividing the three weights by their common factor, the pair consisting of the weighted projective space $\mathbb{P}(12,12,2n+1,4n-4)$ and its branch divisor is isomorphic to the pair $(\mathbb{P},B)$, where
$
\mathbb{P}:=\mathbb{P}(3,3,2n+1,n-1)$
and $B$ is the surface defined by
\[
x_0^{n}+x_1^{n}+x_3x_2+x_1x_3^3=0
\]
in $\mathbb{P}$.
Consequently, the K-polystability of $Y$ is equivalent to that of the log Fano pair
\[
\left(\mathbb{P}, \frac{1}{2}B+\frac{3}{4}H\right),
\]
where $H\subset \mathbb{P}$ is the hypersurface defined by $x_2=0$.



Let $A$ denote the divisor class corresponding to $\mathcal{O}_{\mathbb{P}}(1)$. Define
$\Delta_{\mathbb{P}}:=\frac{1}{2}B+\frac{3}{4}H$ and~$L:=-\left(K_{\mathbb{P}}+\Delta_{\mathbb{P}}\right)$.
Then
$
L\sim_{\mathbb{Q}}\frac{21}{4}A$.

Let $\Pi$  be the hypersurface of $\mathbb{P}$ defined by $x_0=0$. We identify $\Pi$ with the weighted projective plane $\mathbb{P}(3,2n+1,n-1)$ and set
\[
A_{\Pi}:=A|_{\Pi}, \qquad
\ell_B:=B|_{\Pi}, \qquad
\ell_H:=H|_{\Pi}.
\]
The curves $\ell_B$ and $\ell_H$ are defined by
\[
x_1^{n}+x_3x_2+x_1x_3^3=0
\qquad\text{and}\qquad
x_2=0,
\]
respectively. The surface $\Pi$ has three quotient singular points at
\[
o_1=[0:1:0:0], \qquad
o_2=[0:0:1:0], \qquad
o_3=[0:0:0:1]
\]
in $\mathbb{P}$.
By the adjunction formula,
\[
\left(K_{\mathbb{P}}+\frac{1}{2}B+\frac{3}{4}H+\Pi\right)\Big|_{\Pi}
=
K_{\Pi}+\frac{1}{2}\ell_B+\frac{3}{4}\ell_H.
\]
Furthermore,
\begin{equation}\label{eq:G2-LH}
\left(K_{\Pi}+\frac{1}{2}\ell_B+\ell_H\right)\Big|_{\ell_H}
=
K_{\ell_H}
+\frac{2}{3}o_1
+\frac{2n-3}{2n-2}o_3
+\frac{1}{2}o_B,
\end{equation}
where $o_B$ is the unique smooth intersection point of $\ell_B$ and $\ell_H$. It follows that the pair
$
\left(\Pi,\frac{1}{2}\ell_B+\ell_H\right)
$
is purely log terminal along the curve $\ell_H$. Since $\ell_B$ is quasi-smooth, this implies that 
$
\left(\Pi,\frac{1}{2}\ell_B+\frac{3}{4}\ell_H\right)
$
is Kawamata log terminal, and therefore
$
(\mathbb{P},\Delta_{\mathbb{P}}+\Pi)
$
is purely log terminal along $\Pi$.
Moreover, we compute
\[
S_L(\Pi)
=
\frac{1}{L^3}
\int_{0}^{\infty}
\operatorname{vol}_{\mathbb{P}}(L-u\Pi)\,du
=
\int_{0}^{\frac{7}{4}}
\left(1-\frac{4}{7}u\right)^3\,du
=
\frac{7}{16}.
\]



Let  $p$ be a point on $\Pi$, and let $\mathfrak{c}\in |mA_\Pi|$ be an irreducible curve passing through $p$. We have
\begin{equation}\label{eq:G2-2dot}
\begin{split}
S_L(V_{\bullet,\bullet}^{\Pi};\mathfrak{c})
&=
\frac{3}{L^3}
\int_0^{\frac{7}{4}}
\int_0^{\infty}
\operatorname{vol}_{\Pi}\left((L-u\Pi)|_{\Pi}-v\mathfrak{c}\right)\,dv\,du \\
&=
\frac{3}{L^3}
\int_0^{\frac{7}{4}}
\int_0^{\frac{21-12u}{4m}}
\left(\frac{21}{4}-3u-mv\right)^2A_\Pi^2\,dv\,du \\
&=
\frac{21}{16m},
\end{split}
\end{equation}
and
\begin{equation}\label{eq:G2-3dot}
\begin{split}
S_L(W_{\bullet,\bullet,\bullet}^{\Pi,\mathfrak{c}};p)
&=
\frac{3}{L^3}
\int_0^{\frac{7}{4}}
\int_0^{\frac{21-12u}{4m}}
\left(\left((L-u\Pi)|_{\Pi}-v\mathfrak{c}\right)\cdot\mathfrak{c}\right)^2\,dv\,du \\
&=
A_\Pi^2m^2
S_{(\mathbb{P}, \Delta_\mathbb{P})}(V_{\bullet,\bullet}^{\Pi};\mathfrak{c}) \\
&=
\frac{7m}{16(2n+1)(n-1)}.
\end{split}
\end{equation}

Set
$
\Delta_{\Pi}:=\frac{1}{2}\ell_B+\frac{3}{4}\ell_H$.
Let $\ell$ be a curve in the pencil $|(2n+1)A_\Pi|$ passing through $p$.

First,  suppose  that $\ell$ is defined by the equation
$
x_1x_3^2=\alpha x_2
$
for some nonzero constant~$\alpha$. We then have
\[
\left(K_\Pi+\Delta_\Pi+\ell\right)\big|_\ell=
\begin{cases}
K_\ell+\dfrac{7}{6}o_1+\dfrac{4n-3}{4n-4}o_3+\dfrac{1}{2}p_B,
& \text{if }\alpha\neq -1, \\[1ex]
K_\ell+\dfrac{7}{6}o_1+\dfrac{6n-5}{4n-4}o_3,
& \text{if }\alpha=-1,
\end{cases}
\]
where $p_B$ is the unique smooth intersection point of $\ell$ and $\ell_B$ in the case $\alpha\neq -1$. 
Consequently, the pair $(\Pi,\Delta_\Pi+\ell)$ is purely log terminal in a neighborhood of the point $p$ unless $p$ is a singular point.
Applying \cite[Corollary~4.8]{Fu23} together with \eqref{eq:G2-2dot} and \eqref{eq:G2-3dot}, we obtain
\[
\delta_{p}\left(\mathbb{P},\Delta_\mathbb{P}\right)
\ge
\min\left\{
\frac{16}{7},
\frac{16(2n+1)}{21},
\frac{8(n-1)}{7}
\right\}
\ge
\frac{8}{7}
\]
for every point $p\in \ell\setminus\{o_1, o_3\}$.



Next, suppose that $\ell=\ell_H$. The log discrepancy of $\ell_H$ with respect to $(\Pi,\Delta_\Pi)$ is~$\frac{1}{4}$, and for each point $p$ in $\ell_H$ the log 
discrepancy of $p$ with respect to $(\ell_H,\Delta_{\ell_H})$ is at least~$\frac{1}{2(n-1)}$, where $\Delta_{\ell_H}=\frac{2}{3}o_1
+\frac{2n-3}{2n-2}o_3
+\frac{1}{2}o_B$ as in \eqref{eq:G2-LH}.
Then, \cite[Corollary~4.8]{Fu23}, together with \eqref{eq:G2-2dot}, and \eqref{eq:G2-3dot}, immediately yields
\[
\delta_{p}\left(\mathbb{P},\Delta_\mathbb{P}\right)
\ge
\min\left\{
\frac{16}{7},
\frac{4(2n+1)}{21},
\frac{8}{7}
\right\}
\ge
\frac{8}{7}
\]
for every point $p\in\ell_H$. In particular,  $\delta_{o_1}\left(\mathbb{P},\Delta_\mathbb{P}\right)\geq \frac{8}{7}$ and $ \delta_{o_3}\left(\mathbb{P},\Delta_\mathbb{P}\right)\geq \frac{8}{7}$.


Now consider the curve $\ell_3\in |(n-1)A_\Pi|$ defined by $x_3=0$, and suppose that the  point $p$ lies on $\ell_3$. We have
\[
\left(K_\Pi+\Delta_\Pi+\ell_3\right)\big|_{\ell_3}
=
K_{\ell_3}
+\frac{11}{12}o_1
+\frac{5n}{4n+2}o_2.
\]
Again, \eqref{eq:G2-2dot} and \eqref{eq:G2-3dot} imply that
\[
\delta_{p}\left(\mathbb{P},\Delta_\mathbb{P}\right)
\ge
\min\left\{
\frac{16}{7},
\frac{16(n-1)}{21},
\frac{4(2n+1)}{21}
\right\}
\ge
\frac{8}{7}
\]
for every $p\in\ell_3\setminus\{o_2\}$. 


Next, consider the curve $\ell_1\in |3A_\Pi|$ defined by $x_1=0$, and suppose that $p\in \ell_1$. We have
\[
\left(K_\Pi+\Delta_\Pi+\ell_1\right)\big|_{\ell_1}
=
K_{\ell_1}
+\frac{4n+1}{4n+2}o_2
+\frac{4n-3}{4n-4}o_3.
\]
It follows from \eqref{eq:G2-2dot} and \eqref{eq:G2-3dot} that
\[
\delta_{p}\left(\mathbb{P},\Delta_\mathbb{P}\right)
\ge
\min\left\{
\frac{16}{7},
\frac{8(n-1)}{21}
\right\}
\ge
\frac{8}{7}
\]
for every point $p\in \ell_1\setminus\{o_3\}$. 

Consequently,  we have verified that
$\delta_{p}\left(\mathbb{P}, \Delta_\Pi\right)\geq\frac{8}{7}$ for every point $p$ in $\Pi$.



Now consider the point $o_0:=[1:0:0:0]$ in $\mathbb{P}$. The surface $H$ is isomorphic to $\mathbb{P}(1,1,n-1)$, and we identify $H$ with $\mathbb{P}(1,1,n-1)$. Let $A_H$ denote the divisor class on~$H$ corresponding to $\mathcal{O}_H(1)$. Note that
\[
A|_H=\frac{1}{3}A_H.
\]
We have
\[
\left(K_\mathbb{P}+\frac{1}{2}B+H\right)\Big|_{H}
=K_H+\frac{2}{3}\ell_3'+\frac{1}{2}\ell_B',
\]
where $\ell_3'$ is the curve on $H$ defined by $x_3=0$, and $\ell_B':=B|_H$. The curve $\ell_B'$ is defined by
$
x_0^n+x_1^n+x_1x_3=0
$
in $H=\mathbb{P}(1,1,n-1)$.

The curves $\ell_3'$ and $\ell_B'$ intersect at $n$ distinct smooth points
$
p_1,\ldots,p_n
$
of $H$, all of which are different from $o_0$. Note that the pair $(\mathbb{P},\frac{1}{2}B+H)$ is purely log terminal, and the log discrepancy of $H$ with respect to $(\mathbb{P},\Delta_\mathbb{P})$ is $\frac{1}{4}$.

Since
\[
\left(K_H+\frac{1}{2}\ell_B'+\ell_3'\right)\Big|_{\ell_3'}
=
K_{\ell_3'}+\frac{1}{2}\left(p_1+\cdots+p_n\right),
\]
the pair $(H,\frac{1}{2}\ell_B'+\ell_3')$ is purely log terminal. We obtain
\begin{equation}\label{eq:Extra-G2-1dot}
S_L(H)
=
\frac{1}{L^3}\int_0^\infty
\operatorname{vol}_{\mathbb P}(L-uH)\,du
=
\int_0^{\frac{21}{4(2n+1)}}
\left(1-\frac{4(2n+1)}{21}u\right)^3\,du
=
\frac{21}{16(2n+1)},
\end{equation}
\begin{equation}\label{eq:Extra-G2-2dot}
\begin{split}
S_L(V_{\bullet,\bullet}^{H};\ell_3')
&=
\frac{3}{L^3}
\int_0^{\frac{21}{4(2n+1)}}
\int_0^\infty
\operatorname{vol}_{H}
\left(\left.(L-uH)\right|_H-v\ell_3'\right)\,dv\,du\\
&=
\frac{3}{L^3}
\int_0^{\frac{21}{4(2n+1)}}
\int_0^{\frac{21-4(2n+1)u}{12(n-1)}}
\operatorname{vol}_{H}
\left(
\left(\frac{21}{4}-(2n+1)u\right)A|_H-(n-1)vA_H
\right)\,dv\,du\\
&=
\frac{3}{L^3}
\int_0^{\frac{21}{4(2n+1)}}
\int_0^{\frac{21-4(2n+1)u}{12(n-1)}}
\left(
\frac{7}{4}-\frac{2n+1}{3}u-(n-1)v
\right)^2 A_H^2\,dv\,du\\
&=
\frac{7}{16(n-1)},
\end{split}
\end{equation}
and 
\begin{equation}\label{eq:Extra-G2-3dot}
\begin{split}
S_L(W_{\bullet,\bullet,\bullet}^{H,\ell_3'};o_0)
&=
\frac{3}{L^3}
\int_0^{\frac{21}{4(2n+1)}}
\int_0^{\frac{21-4(2n+1)u}{12(n-1)}}
\left(
\left(
\left.(L-uH)\right|_H-v\ell_3'
\right)\cdot\ell_3'
\right)^2\,dv\,du\\
&=
A_H^2(n-1)^2
S_L(V_{\bullet,\bullet}^{H};\ell_3')\\
&=
\frac{7}{16}.
\end{split}
\end{equation}

By \cite[Corollary~4.8]{Fu23}, the values \eqref{eq:Extra-G2-1dot}, \eqref{eq:Extra-G2-2dot}, and \eqref{eq:Extra-G2-3dot} yield
\[
\delta_{o_0}(\mathbb{P},\Delta_\mathbb{P})
\geq
\min\left\{
\frac{4(2n+1)}{21},
\frac{16(n-1)}{21},
\frac{16}{7}
\right\}
\geq\frac{8}{7}.
\]


The automorphism group of the pair $\left(\mathbb{P},\frac12B+\frac34H\right)$ contains the subgroup $\mathfrak{G}$ generated by
\[
\begin{aligned}
[x_0:x_1:x_2:x_3]
&\longmapsto
[\mu_nx_0:x_1:x_2:x_3],\\
[x_0:x_1:x_2:x_3]
&\longmapsto
[x_0:\mu_nx_1:\mu_{3n}x_2:\mu_{3n}^{-1}x_3],
\end{aligned}
\]
where $\mu_r$ denotes a primitive $r$-th root of unity.

Every $\mathfrak{G}$-invariant point other than $o_0$ is contained in the surface $\Pi$. Moreover, every irreducible $\mathfrak{G}$-invariant subvariety of positive dimension intersects $\Pi$. Therefore, the inequality
$
\delta_p\left(\mathbb{P},\Delta_\mathbb{P}\right)\geq\frac{8}{7}
$
for every point $p\in\Pi\cup\{o_0\}$ implies that \[
\delta_\mathfrak{G}\left(\mathbb{P},\Delta_\mathbb{P}\right)\geq\frac{8}{7}>1.
\]
Therefore, by Theorem~\ref{theorem:G-delta}, the pair $(\mathbb{P},\Delta_\mathbb{P})$ is K-polystable, and consequently so is $Y$.



When  $n\equiv 1$  (mod $3$), i.e., $n=3k+1$ for some nonnegative integer $k$, we replace $\mathbb{P}$ at the beginning of the proof by $\mathbb{P}(1,1,2k+1,k)$. The argument then proceeds verbatim.
\end{proof}

\begin{thebibliography}{99}
\bibitem[Fu23]{Fu23} Kento Fujita. \newblock On {K}-stability for {F}ano threefolds of rank 3 and degree 28. \newblock {\em Int. Math. Res. Not. IMRN}, (15):12601--12784, 2023.
\bibitem[Z21]{Z21} Ziquan Zhuang. \newblock Optimal destabilizing centers and equivariant {K}-stability. \newblock {\em Invent. Math.}, 226(1):195--223, 2021.
\end{thebibliography}
\end{document}
